\section{Introduction}
Algorithm Selection is a critical component in the efficient execution of Constraint Programming models, as solver performance varies significantly across different problem instances. Traditional feature extraction methodologies, such as \texttt{fzn2feat}, primarily rely on instance-level statistics that often fail to capture the underlying latent structure and the complex interconnections between variables. While graph-based representations are designed to capture these structural nuances, they are typically coupled with deep neural networks that require substantial computational overhead. This creates a representational gap: non-neural machine learning techniques remain largely ``structure-blind,'' while structure-aware models are often too expensive for practical deployment.

This paper proposes a novel methodology to bridge this gap by implementing a graph conversion algorithm compatible with both classical and deep machine learning paradigms. Our approach utilizes Weisfeiler-Lehman graph kernels to generate robust representations of problem instances. The primary innovation is the introduction of a cut-based representation, which explicitly models structural partitions to provide a more nuanced predictive signal than standard graph features.

To evaluate the efficacy of these features, we constructed three distinct algorithm selection tasks using instances from the 2023, 2024, and 2025 MiniZinc Challenges. Experimental results across Support Vector Machines, Random Forests, and Multi-Layer Perceptrons demonstrate that our features outperform \texttt{fzn2feat} on two out of three tasks, often. Notably, the \texttt{WLc} features often achieved statistically significant improvements ($p < 0.05$), while never performing statistically worse than \texttt{fzn2feat} in any configuration. Our work effectively demonstrates that graph-based structural information can be leveraged to enhance traditional machine learning architectures, ensuring performance gains that are robust across diverse algorithmic structures.